\documentclass[../../book.tex]{subfiles}

\begin{document}

\subsection{Basic Definition}

Polynomials are intensively used in almost all areas of cryptography. In our
particular case, polynomials will encode the information about statements we
will need to prove. That being said, let us define what a polynomial is.

\begin{definition}
    A \textbf{polynomial} $p(T)$ is a function of the form
    \begin{xequation*}
        p(T) = c_0 + c_1 T + c_2 T^2 + \cdots + c_n T^n = \sum_{k=0}^{n} c_k T^k,
    \end{xequation*}
    where $c_0, c_1, \dots, c_n$ are called coefficients of the polynomial.
\end{definition}

Notice that for now we did not specify what are $c_i$'s. We are interested in the case where $c_i \in \mathbb{F}$, where $\mathbb{F}$ is a field. 

\begin{definition}
    A set of polynomials depending on $T$ with coefficients in a field $\mathbb{F}$ is denoted as $\mathbb{F}[T]$, that is
    \begin{xequation}
        \mathbb{F}[T] = \left\{p(T) = \sum_{k=0}^{n} c_k T^k: c_k \in \mathbb{F}, \; k = 0,\dots,n\right\}.
    \end{xequation}
\end{definition}

\begin{example}
    Consider the finite field $\mathbb{F}_3$. Then, some examples of polynomials from $\mathbb{F}_3[T]$ are listed below:
    \begin{enumerate}
        \item $p(T) = 1 + T + 2T^2$.
        \item $q(T) = 1 + T^2 + T^3$.
        \item $r(T) = 2T^3$.
    \end{enumerate}

   If we were to evaluate these polynomials at $1 \in \mathbb{F}_3$, we would get:
    \begin{enumerate}
        \item $p(1) = 1 + 1 + 2 \cdot 1 \; \text{mod} \; 3 = 1$.
        \item $q(1) = 1 + 1 + 1 \; \text{mod} \; 3 = 0$.
        \item $r(1) = 2 \cdot 1 = 2$.
    \end{enumerate}
\end{example}

\vspace{-4mm}

\begin{definition}
    The \textbf{degree} of a polynomial $p(T) = c_0+c_1T+c_2T^2+\dots$ is the largest $k \in \mathbb{Z}_{\geq 0}$ such that $c_k \neq 0$. We denote the degree of a polynomial as $\deg p$. We also denote by $\mathbb{F}[T]^{\leq m} := \{p \in \mathbb{F}[T]: \deg p \leq m\}$.
\end{definition}

\begin{example}
    The degree of the polynomial $p(T) = 1 + 2T + 3T^2$ is $2$. Consequently, $p \in
    \mathbb{F}_5[T]^{\leq 2}$.
\end{example}

\begin{theorem}
    For any polynomials $p,q \in \mathbb{F}[T]$ and $n := \deg p$, $m := \deg q$, the following two statements are true:
    \begin{enumerate}
        \item $\deg (pq) = n + m$.
        \item $\deg (p + q) = \max\{n,m\}$ if $n \neq m$ and $\deg (p+q) \leq m$ for $m=n$.
    \end{enumerate}
\end{theorem}

\subsection{Roots and divisibility}

\begin{definition}
    Let $p \in \mathbb{F}[T]$ be a polynomial of degree $\deg p \geq 1$. A field element $\alpha \in \mathbb{F}$ is called a \textbf{root} of $p$ if $p(\alpha) = 0$.
\end{definition}

\vspace{-4mm}

\begin{example}
    Consider the polynomial $p(T) = 1 + T + T^2 \in \mathbb{F}_3[T]$. Then, $\alpha=1$ is a root of $p(T)$ since $p(\alpha) = 1 + 1 + 1 \; \text{mod} \; 3 = 0$.
\end{example}

One of the fundamental theorems of polynomials is the following.

\begin{theorem}\label{th:root-div}
    Let $p \in \mathbb{F}[T], \deg p \geq 1$. Then, $\alpha \in \mathbb{F}$ is a root of $p$ if and only if there exists a polynomial $q$ (with $\deg q=n-1$) such that
    \begin{equation*}
        p(T) = (T-\alpha)q(T).
    \end{equation*}
\end{theorem}

\begin{example}
    For example, $\alpha=1$ is a root of $p(T) = T^2+2$ over $\mathbb{F}_3$. Indeed, $p(T) = (T-1)(T-2)$ (thus $q(T) = T - 2$ according to Theorem~\ref{th:root-div}).
\end{example}

Also, this might not be immediately obvious, but we can also divide polynomials in the same way as we divide integers. The result of division is not always a polynomial, so we also get a remainder.

\begin{theorem}
    Given $f,g \in \mathbb{F}[T]$ with $g \neq 0$, there are unique polynomials $p,q \in \mathbb{F}[T]$ such that 
    \begin{xequation}
        f = q \cdot g + r, \; 0 \leq \deg r < \deg g.
    \end{xequation}
\end{theorem}

\begin{example}
    Consider $f(T) = T^3+2$ and $g(T) = T+1$ over $\mathbb{R}$. Then, we can write $f(T) = (T^2-T+1)g(T) + 1$, so the remainder of the division is $1$. Typically, we denote this as:
    \begin{xequation*}
        f \, \text{div} \, g = T^2-T+1, \quad f \, \text{mod} \, g = 1.
    \end{xequation*}

    The notation is similar to the one used in integer division.
\end{example}

Similarly, one can define $\text{gcd}$, $\text{lcm}$, and other number field theory operations for polynomials\footnote{This follows from the fact that both $\mathbb{Z}$ and $\mathbb{F}[T]$ are so-called \textit{principal ideal domains}. For a careful treatment of the details, see for example \cite[Chapters 1 \& 2]{ireland}.}. However, we will not go into details here, besides mentioning the divisibility.

\begin{definition}
    A polynomial $f \in \mathbb{F}[T]$ is called \textbf{divisible} by $g \in \mathbb{F}[T]$ (or, $g$ \textbf{divides} $f$, written as $g \mid f$) if there exists $h \in \mathbb{F}[T]$ such that $f=gh$.
\end{definition}

\begin{proposition}[Theorem~\ref{th:root-div} restated]
    If $\alpha \in \mathbb{F}$ is a root of polynomial $p \in \mathbb{F}[T]$, then $(T-\alpha) \mid p$.
\end{proposition}

\begin{definition}
    A polynomial $f \in \mathbb{F}[T]$ is said to be \textbf{irreducible} in $\mathbb{F}$ if there are no polynomials $g,h \in \mathbb{F}[T]$ both of positive degree such that $f = gh$.
\end{definition}

\begin{example}
    A polynomial $f(T) = T^2+16$ is irreducible in $\mathbb{R}$. On the other hand, $f(T) = T^2-2$ is not irreducible since $f(T) = (T-\sqrt{2})(T+\sqrt{2})$. 
\end{example}

\begin{example}
    There are no polynomials over complex numbers $\mathbb{C}$ with degree greater than $1$ that are irreducible. This follows from the \textit{fundamental theorem of algebra}.
\end{example}

\subsection{Interpolation}

Now, let us ask the question: what defines the polynomial? 
Given expression $p = \sum_{k=0}^n c_kT^k$, one can respond: ``I just need to know the coefficients $\{c_k\}_{k=0}^n$''.

Indeed, each polynomial of degree $n$ is uniquely determined by the vector of its coefficients $(c_0,c_1,\dots,c_n) \in \mathbb{F}^{n+1}$. However, that is not the only way to define a polynomial.

Suppose that $p(T) = aT + b$ -- just a simple linear function over $\mathbb{R}$. Suppose that $p$ intercepts $(0,0)$ and $(1,2)$. Then, you can easily say that $p(T) = 2T$. 

The more general question is: suppose $\deg p = n$, how many points do I need to define the polynomial $p(T)$ uniquely? 
The answer is $n+1$ distinct points. 
This is the idea behind the interpolation:
the polynomial is uniquely defined by $n+1$ distinct points on the plane. An example is depicted 
in \Cref{fig:interpolation}. Now, let us see how we can interpolate the polynomial practically.

\begin{figure}[H]
    \centering
    \includegraphics[width=0.8\textwidth]{contents/images/lecture_1/interpolation.pdf}
    \caption{$5$ points on the plane uniquely define the polynomial of degree $4$.}
    \label{fig:interpolation}
\end{figure}

\begin{theorem}
    Given a set of points $\{(x_0,y_0),(x_1,y_1),\dots,(x_n,y_n)\} \subseteq
    \mathbb{F} \times \mathbb{F}$, there is a unique polynomial $L$ of degree
    $n$ such that $L(x_i) = y_i$ for all $i=0,\dots,n$. This polynomial is
    called the \textbf{Lagrange interpolation polynomial} and can be found
    through the following formula:
    \begin{xequation}
        L(T) = \sum_{i=0}^{n} y_i \ell_i(T), \quad \ell_i(T) = \prod_{j=0, j \neq i}^{n} \frac{T-x_j}{x_i-x_j}.
    \end{xequation}
\end{theorem}

\begin{lemma}
    The polynomials $\{\ell_i\}_{i=0}^n$, in fact, have quite an interesting property:
    \begin{xequation*}
        \ell_i(x_j) = \delta_{ij} = \begin{cases}
            1, & i = j \\ 0, & i \neq j
        \end{cases},
    \end{xequation*}

    where $\delta_{ij}$ is the so-called \textit{Kronecker delta}. Moreover, $\{\ell_i\}_{i=0}^n$ form a \textit{basis} of $\mathbb{F}[T]^{\leq n}$ (see \Cref{section:linear-algebra} for the definition). Put simply, for any polynomial $p \in \mathbb{F}[T]^{\leq n}$ there exist unique coefficients $\alpha_0,\dots,\alpha_n \in \mathbb{F}$ such that
    \begin{xequation}
        p(T) = \sum_{i=0}^{n} \alpha_i \ell_i(T).
    \end{xequation}
\end{lemma}

\begin{example}
    Suppose we have points $(0,1)$ and $(1,2)$. Then, the Lagrange interpolation polynomial is
    \begin{equation*}
        L(T) = 1 \cdot \frac{T-1}{0-1} + 2 \cdot \frac{T-0}{1-0} = (-1) \cdot (T-1) + 2 \cdot T = T + 1.
    \end{equation*}
\end{example}

\subsection{Schwartz-Zippel Lemma}

Schwartz-Zippel Lemma is central for most modern zero-knowledge protocols. Its
simple interpretation is the following: if we want to check whether two
polynomials are equal, we can do it by checking whether they are equal at a
random point. The probability of two different polynomials being equal at a
random point is very low for sufficiently large finite fields $\mathbb{F}$.
However, let us formulate the lemma more rigorously. Let us start with the
single-variable case, which is easiest to understand.

\vspace{-0.6mm}

\begin{lemma}[Single Variable Schwartz-Zippel Lemma]\label{lemma:one-sz} Let
    $\mathbb{F}$ be a finite field and let $f \in \mathbb{F}[T]$ be a non-zero
    polynomial. Then,
    \begin{xequation}
        \underset{\alpha \leftarrowS \mathbb{F}}{\text{Pr}}
        [f(\alpha) = 0] \leq \frac{\deg f}{\#\mathbb{F}}.
    \end{xequation}
\end{lemma}

\vspace{-0.6mm}

\textbf{Proof Idea.} The proof is quite simple. The polynomial of degree up to
$n:=\deg f$ can have up to $n$ roots. Therefore, when we pick the point $\alpha$
uniformly at random from the field $\mathbb{F}$, there is at most $n/\#\mathbb{F}$
chance to pick the root of the polynomial.

However, the lemma can be generalized to the multi-variable case, which is more
difficult to prove. The generalization is as follows.

\vspace{-0.6mm}

\begin{lemma}[Multivariable Schwartz-Zippel] \label{lemma:sz} Let $\mathbb{F}$
    be a finite field and $S \subseteq \mathbb{F}^n$ be a finite
    $n$-dimensional subspace. Let $f: S \to \mathbb{F}$ be a polynomial
    in $n$ variables. Then,
    
    \vspace{-1.1mm}
    
    \begin{xequation}
        \underset{\bm{\alpha} \leftarrowS S}{\text{Pr}}
        [f(\bm{\alpha}) = 0] \leq \frac{\deg f}{\#S}.
    \end{xequation}
\end{lemma}

\vspace{-4.1mm}

\begin{example}
Suppose we are working over the Mersenne prime field of order $p=2^{31}-1$ with
a polynomial $f(T,W) = T^{50} + 3TW + W^5$. If we pick $(\alpha,\beta)$
at random from $\mathbb{F}_p^2$, the probability of $f(\alpha,\beta) = 0$ is at most
$50/(2^{31}-1)$, providing approximately $25$ bits of security.
\end{example}

\subsection{Some Fun: Shamir's Secret Sharing}

Shamir's Secret Sharing, also known as $(t,n)$-threshold scheme, is one of the protocols exploiting Lagrange Interpolation. 

But first, let us define what secret sharing is. Suppose we have a secret data $\alpha$, which is represented as an element from some finite set $F$. We divide this secret into $n$ pieces $\alpha_1,\dots,\alpha_n \in F$ in such a way:
\begin{enumerate}
    \item Knowledge of any $t$ shares can reconstruct the secret $\alpha$.
    \item Knowledge of any number of shares below $t$ cannot be used to reconstruct the secret $\alpha$.
\end{enumerate}

Now, let us define the sharing scheme.

\begin{definition}
    \textbf{Secret Sharing} scheme is a pair of efficient algorithms $(\mathsf{Gen}, \mathsf{Comb})$ which work as follows:
    \begin{itemize}
        \item $\mathsf{Gen}(\alpha, t, n)$: probabilistic sharing algorithm that yields $n$ shards $(\alpha_1,\dots,\alpha_t)$ for which $t$ shards are needed to reconstruct the secret~$\alpha$.
        \item $\mathsf{Comb}(\mathcal{I}, \{\alpha_i\}_{i \in \mathcal{I}})$: deterministic reconstruction algorithm that reconstructs the secret $\alpha$ from the shards $\mathcal{I} \subset \{1,\dots,n\}$ of size~$t$.
    \end{itemize}

    Here, we require the \textbf{correctness}: for every $\alpha \in F$, for every possible output $(\alpha_1,\dots,\alpha_n) \gets \mathsf{Gen}(\alpha, t, n)$, and any $t$-size subset $\mathcal{I}$ of $\{1,\dots,n\}$ we have
    \begin{xequation}
        \mathsf{Comb}(\mathcal{I}, \{\alpha_i\}_{i \in \mathcal{I}}) = \alpha.
    \end{xequation}
\end{definition}

Now, Shamir's protocol is one of the most famous secret sharing schemes. It works as follows: our finite set is $\mathbb{F}_q$ for some large prime $q$. Then, algorithms in the protocol are defined as follows:
\begin{itemize}
    \item $\mathsf{Gen}(\alpha, k, n)$: choose random $k_1,\dots,k_{t-1}
    \leftarrowS \mathbb{F}_q$ and define polynomial $\omega(T) := \alpha +
    \sum_{j=1}^{t-1}k_jT^j \in \mathbb{F}_q[T]^{<t}$. Then, compute
    $\alpha_i \gets \omega(i) \in \mathbb{F}_q$ for each $i$ and return
    $(\alpha_1,\dots,\alpha_n)$.
    \item $\mathsf{Comb}(\mathcal{I}, \{\alpha_i\}_{i \in \mathcal{I}})$: reconstruct the polynomial $\omega$ using Lagrange interpolation and return $\omega(0) = \alpha$.
\end{itemize}

The combination function is possible since, having $t$ points $\{(i,\alpha_i)\}_{i
\in \mathcal{I}}$ with $\omega(i) = \alpha_i$, we can fully reconstruct the
polynomial $\omega$ and then evaluate it at $0$ to get $\alpha$. Instead,
suppose we have only $t-1$ (or less) pairs $\{(i,\alpha_i)\}_{i \in \mathcal{I}}$.
Then, there are many polynomials $\omega$ that pass through these points (in
fact, if we were in the field of real numbers, this number would be infinite),
and thus the secret $\alpha$ is not uniquely determined. The intuition behind
Shamir's Secret Sharing protocol is illustrated in Figure~\ref{fig:shamir}: for
$t=3$, having only two \textcolor{blue}{blue} points, we cannot reconstruct the
polynomial, while knowing the third \textcolor{red}{red} point allows us to
uniquely determine the polynomial and thus get value at $0$.

\vspace{-2mm}

\begin{figure}[H]
    \centering
    \includegraphics[width=0.5\textwidth]{contents/images/lecture_1/shamir_demo.pdf}
    \caption{Illustration of Shamir's Secret Sharing}
    \label{fig:shamir}
\end{figure}

\subsubsection{SageMath Implementation}\label{sec:shamir-impl}

Let us implement this protocol in SageMath.
First, we will set up our finite field $\mathbb{F}_q$ and polynomial ring $\mathbb{F}_q[T]$.

\begin{minted}{python}
from typing import List, Tuple

q = 2^(31) - 1
Fq = GF(q)
P.<x> = PolynomialRing(Fq)
\end{minted}

Now, let us write the $\mathsf{Gen}$ function that returns shares $\{(i, f(i))\}_{i \in [n]}$.

\begin{minted}{python}
def generate(n: int, k: int) -> List[[Fq, int]]:
    """
    Generate a Shamir secret sharing scheme with n shares and k threshold.
    """

    # Generate a random polynomial
    f = P.random_element(degree=k-1)
    
    # Generate a secret
    secret = Fq.random_element()
    print(f'Chosen secret: {secret}')

    # Set the constant term as the secret
    f = P([secret] + [f[i] for i in range(1, k)])
    print(f'Polynomial: {f}')
    
    # Return shares
    return [(Fq(i), f(i)) for i in range(1, n+1)]
\end{minted}

This way, using the following function, we generate shares:

\vspace{-6px}
\begin{minted}{python}
shares = generate(5, 3) # Generate 3-of-5 shares
\end{minted}
\vspace{-6px}

Finally, we write the $\mathsf{Comb}$ function that reconstructs the secret value by interpolating given values.

\begin{minted}{python}
def reconstruct(shares: List[Tuple[int, Fq]], threshold: int) -> Fq:
    """
    Reconstruct the secret from the shares using Lagrange interpolation.
    """
    
    assert len(shares) >= threshold, f"Not enough shares to reconstruct 
                                    the secret, got {len(shares)}"

    reconstructed_polynomial = P.lagrange_polynomial(shares)
    print(f'Reconstructed polynomial: {reconstructed_polynomial}')

    # The secret is the constant term of the polynomial
    return reconstructed_polynomial(0)

recovered_secret = reconstruct(shares[1:4], 3) # Try any subset of 3 elements
print(f'Recovered secret: {recovered_secret}')
\end{minted}

\subsection{Exercises}

\begin{xexercise}
    {Exercise 1.}
    {
        Suppose for three polynomials $p,q,r \in \mathbb{F}[T]$ we have $\deg p = 3, \deg q = 4, \deg r = 5$. 
        Which of the following is true for $n := \deg \{(p-q)r\}$?
    }
    {2}
    {
        \item $n = 9$.
        \item $n$ might be less than $9$.
        \item $n = 20$.
        \item $n$ is less than $\deg \{qr\}$. 
    }
\end{xexercise}

\begin{xexercise}
    {Exercise 2.}
    {
        Define the polynomial over $\mathbb{F}_5$: $f(T) := 4T^2 + 2$. Which of the following is the root of $f$?
    }
    {2}
    {
        \item $2$
        \item $3$
        \item $4$
        \item $f$ has no roots.
    }
\end{xexercise}

\begin{xexercise}
    {Exercise 3.}
    {
        Quadratic polynomial $p(T) = aT^2+bT+c \in \mathbb{R}[T]$ has zeros at $1$ and $2$ and $p(0) = 2$. 
        Find the value of $a+b+c$.
    }
    {2}
    {
        \item $0$
        \item $-1$
        \item $1$
        \item Not enough information is provided.
    }
\end{xexercise}

\begin{xexercise}
    {Exercise 4*.}
    {
        Consider two polynomials: 
        \begin{xequation}
            p_n(T) = T^{2n} + T^n -2, \quad q(T) = T^2 + 1
        \end{xequation}
        For which values of $n$ does polynomial $q$ divide $p_n$?
    }
    {1}
    {
        \item For all even $n$.
        \item For all odd $n$.
        \item For all $n$ which are divisible by 4.
        \item Just for $n = 4$.
        \item It does not work for any natural number $n$.
    }
\end{xexercise}

\begin{xexercise}
    {Exercise 5.} { Suppose $f \in \mathbb{F}_p[T]$ is a $d$-degree polynomial
    with $d$ \textbf{distinct} roots in $\mathbb{F}_p$. Suppose Alice wants to
    fool Bob by claiming that $f$ is not a zero polynomial. The protocol
    consists of $n$ rounds, where on each round Bob samples the random scalar
    $\tau \leftarrowS \mathbb{F}_p$ and checks the value of $f(\tau)$. What
    is the probability that Alice successfully fools Bob during the protocol?}
    {2} {
        \item Exactly $(d/p)^n$
        \item Strictly less than $(d/p)^n$
        \item Exactly $nd/p$.
        \item Strictly less than $d/np$.
    }
\end{xexercise}

\begin{xexercise}
    {Exercise 6.} { Suppose $f \in \mathbb{F}_p[T_1,\dots,T_n]$ is an
    $n$-variable polynomial. Carol randomly samples the point $\bm{\alpha}
    \leftarrowS S$ where $S \subseteq \mathbb{F}_p^n$ is a
    set of points $(\alpha_1,\dots,\alpha_n)$ such that $\sum_{i=1}^n \alpha_i=1$ over
    $\mathbb{F}_p$. According to the Schwartz-Zippel Lemma, what is the tightest
    upper bound for $\Pr[f(\bm{\alpha})=0]$ among specified below?}
    {1} {
        \item $\deg f/p^n$
        \item $\deg f/p^{n-1}$
        \item $\deg f/p^{n-2}$
        \item $\deg f/p$
    }
\end{xexercise}

\end{document}
